\section{Method overview}
\label{sec:method}

This section describes the core algorithmic structure studied in this paper
and defines the quantities used in the theoretical statements. \cref{alg:citree}
gives the end-to-end overview. Implementation-aligned stage routines and
additional details are collected in \cref{app:methods}.

\paragraph{Reference mode vs.\ default engineering mode.}
The manuscript distinguishes two execution regimes. The \emph{reference mode}
used by the theorem chain fixes the resampling budget in advance, exhaustively
evaluates the tested Stage~A family at the node of interest, and interprets the
resulting Stage~A p-values only at fixed nodes under the complete permutation
null. The \emph{default engineering mode} used by the software may additionally
enable adaptive early stopping, candidate-order heuristics, subtree feature
muting, threshold subsampling, and forest bootstrap resampling. These options
can improve runtime substantially, but they change the inferential object and
are treated below as algorithmic or empirical rather than as part of the
fixed-$B$ theorem chain.

\subsection{Node expansion: Stage A and Stage B}
\label{sec:stageA-stageB}

Consider a fixed node $t$ with data $(X_t, Y_t)$ and a candidate feature set
$F_t$ of size $m_t = |F_t|$. Let $I_t \subseteq \{1,\dots,n\}$ denote the sample
indices at node $t$ (size $n_t$) and write $X_t := (X_i)_{i \in I_t}$ and
$Y_t := (Y_i)_{i \in I_t}$ (see \cref{app:setup} for a notation summary).
In the recursive algorithm we distinguish: the \emph{available} feature set
$F_{\mathrm{avail}}$ passed down the recursion; the \emph{descendant-safe} subset
$F_{\mathrm{desc}} \subseteq F_{\mathrm{avail}}$ of non-constant features at $t$;
and the \emph{tested} subset $F_t \subseteq F_{\mathrm{desc}}$ after applying
\texttt{max\_features}.
Stage~A returns an updated feature set $F_{\mathrm{avail}}'$ that is passed to
descendants (\cref{alg:citree}); by default $F_{\mathrm{avail}}' =
F_{\mathrm{desc}}$, while optional heuristics may further prune it
(\cref{app:methods-heuristics}).

\paragraph{Stage A (feature screening).}
For each feature $j \in F_t$, the procedure computes a selector statistic
$T^{\mathrm{sel}}_j(X_{t,j}, Y_t)$ (larger indicates stronger association under a
right-tail convention) and a Monte Carlo permutation p-value $p_{t,j}$.
It selects
%
\begin{equation}
  j_t^\star := \operatorname{uargmin}_{j \in F_t} p_{t,j}, \qquad
  p_t^\star := \min_{j \in F_t} p_{t,j},
\end{equation}
%
where $\operatorname{uargmin}$ breaks ties uniformly at random among minimizers, and proceeds to Stage~B only if
$p_t^\star$ is strictly below the configured significance
level $\alpha_{\mathrm{sel}}$ (optionally Bonferroni-adjusted by $m_t$, i.e.,
compare to $\alpha_{\mathrm{sel}}/m_t$).
In the fixed-$B$ reference mode with an exhaustive scan over $F_t$, this corresponds to
evaluating all tested features and selecting the (tie-broken) minimum. In
the default engineering configuration, the implementation may instead order
candidates by a cheap score and stop the scan once a significant feature has
been found, so the selected feature is then the first sufficiently significant
candidate in the scan order rather than necessarily the global minimizer over
all tested features. Sequential permutation testing can also stop drawing
permutations early. Unless stated otherwise, the formal guarantees in
\cref{sec:theory} are limited to fixed-$B$ testing with an exhaustive tested
family at a fixed node.
Because $p_t^\star$ is a minimum over $m_t$ tests, it is not a
selection-adjusted p-value for the selected feature. Under the nodewise
\emph{complete} null at $t$, a conservative complete-null decision-p-value
summary for the Stage~A \emph{decision} is
$\tilde p_t^{\mathrm{sel}} := \min(1,\, m_t p_t^\star)$. If this summary is
thresholded using the same strict comparison, i.e., $\tilde p_t^{\mathrm{sel}} < \alpha_{\mathrm{sel}}$,
it yields the same reject/accept rule as checking whether
$p_t^\star < \alpha_{\mathrm{sel}}/m_t$.
Under the nodewise complete null at $t$ (all tested features are null), this
Bonferroni rule controls the probability of making any Stage~A rejection at level
$\alpha_{\mathrm{sel}}$ in fixed-$B$ mode under the assumptions in
\cref{sec:theory} (see \cref{prop:stageA-global-null}).
In this manuscript, the inferential claim is for the Stage~A decision event at
the fixed node; we do not interpret $p_t^\star$ as a calibrated p-value for the
selected feature.
Implementation-aligned Stage~A pseudocode is given in
\cref{alg:stageA} (\cref{app:methods-stageA}).

\paragraph{Stage B (threshold screening).}
Given $j_t^\star$, the procedure constructs a finite threshold set $C_{t,j_t^\star}$
and evaluates split quality over candidates. When permutation testing is
enabled for Stage~B, this is treated as a left-tail test over candidate
thresholds. In the implementation, the permutation-tested Stage~B statistic is
the \emph{unweighted} sum of child impurities, while deterministic threshold
ordering heuristics and the final \texttt{min\_impurity\_decrease} check use
weighted impurity. Because $c_t^\star$ and the reported minimum p-value are
selected over $C_{t,j_t^\star}$ (and because $j_t^\star$ itself is selected
using the same labels), in this paper we treat Stage~B p-values as
post-selection/adaptive statistics unless additional selective-inference
machinery or sample splitting is used.
Implementation-aligned Stage~B pseudocode and threshold-set details are given in
\cref{alg:stageB} and \cref{app:methods-thresholds}.

\begin{algorithm}[t]
  \caption{Conditional inference tree training (Stage A $\rightarrow$ Stage B)}
  \label{alg:citree}
  \begin{algorithmic}[1]
    \Require Training data $(X,Y)$, hyperparameters $\theta$
    \Ensure Trained tree $T$
    \State $p \gets$ number of columns of $X$
    \State \Return \Call{GrowNode}{$X,Y,1,\{1,\dots,p\},\theta$}
    \Procedure{GrowNode}{$X_t,Y_t,d,F_{\mathrm{avail}},\theta$}
      \If{\Call{Stop}{$X_t,Y_t,d,\theta$}}
        \State \Return \Call{Leaf}{$Y_t$}
      \EndIf
      \State $(j_t^\star, p_t^\star, m_t, F_{\mathrm{avail}}') \gets \Call{StageA}{X_t,Y_t,F_{\mathrm{avail}},\theta}$ \Comment{see \cref{alg:stageA}}
      \If{$j_t^\star = \bot$} \Comment{e.g., no candidates, or $p_t^\star \ge \alpha_{\mathrm{sel}}/m_t$ if Bonferroni}
        \State \Return \Call{Leaf}{$Y_t$}
      \EndIf
      \State $(c_t^\star, p_t^{\mathrm{split}}) \gets \Call{StageB}{X_t,Y_t,j_t^\star,\theta}$ \Comment{see \cref{alg:stageB}}
      \If{$c_t^\star = \bot$}
        \State \Return \Call{Leaf}{$Y_t$}
      \EndIf
      \State $(X_t^L,Y_t^L,X_t^R,Y_t^R) \gets$ \Call{Partition}{$X_t,Y_t,j_t^\star,c_t^\star$}
      \If{\Call{ImpurityDecrease}{$Y_t,Y_t^L,Y_t^R,\theta$} $< \texttt{min\_impurity\_decrease}$}
        \State \Return \Call{Leaf}{$Y_t$}
      \EndIf
      \State $u \gets$ internal node with split $(j_t^\star,c_t^\star)$ and stored $(p_t^\star,p_t^{\mathrm{split}})$
      \State $u.\mathrm{left} \gets \Call{GrowNode}{X_t^L,Y_t^L,d+1,F_{\mathrm{avail}}',\theta}$
      \State $u.\mathrm{right} \gets \Call{GrowNode}{X_t^R,Y_t^R,d+1,F_{\mathrm{avail}}',\theta}$
      \State \Return $u$
    \EndProcedure
  \end{algorithmic}
\end{algorithm}

\subsection{Interpretation of reported p-values (scope summary)}
\label{sec:pvalue-scope}

We use the term ``p-value'' in two different roles: (i)~as a classical
Monte Carlo permutation p-value for fixed-node Stage~A screening in fixed-$B$
mode, and (ii)~as an algorithmic statistic for adaptive parts of tree growth.
Adaptivity enters both through the recursive training procedure (internal nodes
and Stage~B after selecting a feature using the same labels) and through
computational heuristics (e.g., sequential permutation testing and early-terminated scans over candidates).
In particular:
\begin{itemize}
  \item \textbf{Stage A at a fixed node (fixed-$B$, exhaustive scan):} per-feature +1 Monte Carlo
  permutation p-values are super-uniform under the nodewise complete-null
  exchangeability and label-independence assumptions in \cref{sec:theory}. The selected-feature
  minimum $p_t^\star$ is not selection-adjusted; the Bonferroni decision can be
  summarized by the conservative complete-null decision-p-value summary
  $\tilde p_t^{\mathrm{sel}}=\min(1,m_t p_t^\star)$ under the nodewise complete null.
  \item \textbf{Stage B threshold p-values:} these are computed after selecting
  $j_t^\star$ using the same labels, so they are post-selection/adaptive. We
  treat them as split-selection scores unless additional post-selection
  machinery or sample splitting is used.
  \item \textbf{Internal-node p-values:} internal nodes are data-adaptive
  because node membership depends on earlier label-dependent splits. Without
  additional machinery, internal-node p-values should be treated as algorithmic
  statistics rather than calibrated p-values.
  \item \textbf{Early-stopped outputs:} sequential permutation testing changes the
  distribution of the reported running estimate; in this manuscript we do not
  claim optional-stopping validity for early-stopped values.
  \item \textbf{Early-terminated scans:} when candidate scans terminate early (e.g., after finding a
  significant candidate), the selected feature/threshold need not be the global minimizer over all candidates, even if
  each per-candidate p-value is computed with a fixed resampling budget.
\end{itemize}

\begin{table}[t]
  \centering
  \small
  \setlength{\tabcolsep}{5pt}
  \begin{tabular}{@{}
    >{\raggedright\arraybackslash}p{0.22\linewidth}
    >{\raggedright\arraybackslash}p{0.30\linewidth}
    >{\raggedright\arraybackslash}p{0.40\linewidth}
  @{}}
    \toprule
    Quantity & Where it arises & Interpretation in this manuscript \\
    \midrule
    $p_{t,j}$ (+1 MC) &
    Stage~A, fixed node, fixed-$B$ &
    Fixed-node permutation p-value; super-uniform under A0.1--A0.5 (complete-null exchangeability). \\
    $p_t^\star = \min_{j\in F_t} p_{t,j}$ &
    Stage~A feature selection &
    Not selection-adjusted; the nodewise Bonferroni decision can be summarized by $\tilde p_t^{\mathrm{sel}}=\min(1,m_t p_t^\star)$ under the nodewise complete null. \\
    $p_t^{\mathrm{split}}$ &
    Stage~B threshold selection &
    Post-selection/adaptive (computed after selecting $j_t^\star$ using the same labels); treated as a split-selection score unless additional machinery is used. \\
    Internal-node p-values &
    Recursive tree growth &
    Algorithmic statistics (node membership depends on previous label-dependent splits). \\
    Early-stopped $\widehat p$ &
    Sequential permutation testing &
    Algorithmic statistic; no optional-stopping validity claims in this manuscript. \\
    \bottomrule
  \end{tabular}
  \caption{Scope of p-value interpretations used in this manuscript.}
  \label{tab:pvalue-scope}
\end{table}

\subsection{Monte Carlo permutation p-values (+1 correction)}
\label{sec:plusone}

Let $T_0$ be the observed statistic and $T_1,\dots,T_B$ the statistics computed
on $B$ i.i.d.\ random label permutations. For a right-tail test (larger is more
extreme), we use the Phipson--Smyth (+1) Monte Carlo p-value
\citep{NorthCurtisSham2002EmpiricalPValuesMonteCarlo,PhipsonSmyth2010PermutationPValues}
%
\begin{equation}
  p_{\mathrm{MC}}^{\ge} := \frac{1 + \sum_{b=1}^B \ind\{T_b \ge T_0\}}{B+1}.
\end{equation}
%
For a left-tail test (smaller is more extreme), define
%
\begin{equation}
  p_{\mathrm{MC}}^{\le} := \frac{1 + \sum_{b=1}^B \ind\{T_b \le T_0\}}{B+1}.
\end{equation}
%
In \citrees, Stage~A uses the right-tail convention and Stage~B uses the
left-tail convention. In both cases, the implementation uses the conservative
tie rule that counts ties as exceedances.

\subsection{Multi-selector mode (max-statistic)}
\label{sec:multiselector}

In multi-selector mode, multiple association scores are combined within a
feature by taking a maximum across selector statistics inside each permutation.
This yields a single permutation p-value per feature and supports robust
screening when different selectors have complementary power.

\subsection{Forests and bootstrap resampling}
\label{sec:bootstrap}

The forest procedure fits $M$ trees and aggregates their predictions. To induce
diversity across trees, the method supports optional bootstrap resampling of
rows. In the simplest case this is the classical bootstrap, where each tree
draws $n$ indices with replacement from $\{1,\dots,n\}$ with uniform
probabilities $\pi_i = 1/n$ \citep{Efron1979Bootstrap}.

For classification forests, we additionally support class-conditional bootstrap
sampling to control class composition within each bootstrap sample: stratified
sampling performs draws within each class so that expected proportions match
the training data, undersampling draws $n_{\min}$ samples per class (where
$n_{\min}$ is the minority class count), and oversampling draws a fixed total
of $n$ (or a configured bootstrap sample cap) samples with per-class counts as
equal as possible.
These options are engineering heuristics and are not part of this manuscript's
formal permutation-test guarantees. In particular, the root-level theorem chain
in \cref{sec:theory} should be read as applying to trees or forests only when
the root resampling scheme preserves the required permutation exchangeability
assumption.
